\chapter{Comunicación entre servicios}
\label{cap:AnexoC}

Este anexo documenta los protocolos y contratos de comunicación entre los servicios del sistema: la API REST entre frontend y backend, el canal WebSocket para mensajes en tiempo real y la llamada HTTP del backend al motor de IA. La Figura~\ref{fig:conexion-fe-be} ofrece una visión de conjunto.

\begin{figure}[ht]
\centering
\resizebox{\textwidth}{!}{
\begin{tikzpicture}[
  every node/.style={font=\small\sffamily},
  comp/.style={draw=aquaESI!80, fill=aquaESI!15, thick, rectangle, rounded corners=3pt,
               minimum width=3.0cm, minimum height=1.6cm, align=center},
  lbl/.style={font=\small\sffamily, text=black, fill=white, inner sep=1.5pt, align=center},
  arr/.style={->, >=Stealth, thick},
  brr/.style={<->, >={Stealth[length=3mm]}, thick, shorten <=1pt, shorten >=1pt},
]
\node[comp] (fe) at (0,0) {\textbf{Frontend}\\\footnotesize Next.js / React};
\node[comp] (be) at (7.5,0) {\textbf{Backend}\\\footnotesize Phoenix (Elixir)};
\node[comp] (ai) at (14.5,0) {\textbf{AI Engine}\\\footnotesize FastAPI (Python)};

\draw[brr] ([yshift=6mm]fe.east) -- node[lbl, above]{REST: auth, conversaciones,\\ cie10, traducciones} ([yshift=6mm]be.west);
\draw[brr, ocre!85] ([yshift=-6mm]fe.east) -- node[lbl, below]{\color{black}WebSocket: análisis y\\ validación en tiempo real} ([yshift=-6mm]be.west);

\draw[arr, turquesa!85] (be.east) -- node[lbl, above]{HTTP: \texttt{/predict},\\ \texttt{/explain}} (ai.west);
\end{tikzpicture}
}
\caption{Esquema de comunicación entre servicios: el frontend habla con el backend por API REST (autenticación y consultas) y por un canal WebSocket (análisis en tiempo real); el backend invoca al motor de IA por HTTP.}
\label{fig:conexion-fe-be}
\end{figure}

\section{Flujo de análisis clínico}

La Figura~\ref{fig:seq-analisis} ilustra la secuencia completa de mensajes que se intercambian cuando el codificador envía un informe clínico al sistema.

\begin{figure}[ht]
\centering
\resizebox{\textwidth}{!}{-% Diagrama de secuencia, Análisis de informe clínico
% Uso: -% Diagrama de secuencia, Análisis de informe clínico
% Uso: -% Diagrama de secuencia, Análisis de informe clínico
% Uso: \input{figs/seq-analisis} dentro de figure+\centering+\resizebox{\textwidth}{!}{...}
%
% Participantes (x): Usuario=0, Frontend=3.6, Backend=7.2, AI Engine=11
% Tiempo (y): crece hacia abajo (negativo en tikz)

\begin{tikzpicture}[
 every node/.style={font=\small\sffamily},
 % Cabecera de participante
 phead/.style={
 draw=aquaESI!80, fill=aquaESI!20, thick,
 rectangle, rounded corners=3pt,
 minimum width=2.4cm, minimum height=0.75cm, align=center
 },
 % Línea de vida
 life/.style={draw=gray!40, thick, dashed},
 % Mensaje síncrono (flecha llena)
 sync/.style={->, >=Stealth, thick, aquaESI!80},
 % Mensaje de retorno (flecha discontinua)
 ret/.style={->, >=Stealth, dashed, thick, gray!60},
 % Nota lateral
 note/.style={
 draw=ocre!60, fill=ocre!10,
 rectangle, rounded corners=2pt,
 font=\scriptsize\sffamily, inner sep=3pt, align=left
 },
 % Etiqueta de mensaje
 msglbl/.style={font=\scriptsize\sffamily, fill=white,
 inner sep=1pt, above, midway}
]

%% ---- COORDENADAS X de cada participante ----
\def\xu{0} % Usuario
\def\xf{3.6} % Frontend
\def\xb{7.2} % Backend
\def\xa{11.0} % AI Engine

%% ---- CABECERAS ----
\node[phead] at (\xu, 0) {\textbf{Usuario}};
\node[phead] at (\xf, 0) {\textbf{Frontend}};
\node[phead] at (\xb, 0) {\textbf{Backend}};
\node[phead] at (\xa, 0) {\textbf{AI Engine}};

%% ---- LÍNEAS DE VIDA ----
\draw[life] (\xu,-0.38) -- (\xu,-13.0);
\draw[life] (\xf,-0.38) -- (\xf,-13.0);
\draw[life] (\xb,-0.38) -- (\xb,-13.0);
\draw[life] (\xa,-0.38) -- (\xa,-13.0);

%% ---- ACTIVACIONES (barras sobre línea de vida) ----
% Backend activo durante la llamada al AI Engine
\filldraw[draw=aquaESI!70, fill=aquaESI!25]
 (\xb-0.12,-2.2) rectangle (\xb+0.12,-11.6);
% AI Engine activo durante la predicción
\filldraw[draw=turquesa!70, fill=turquesa!20]
 (\xa-0.12,-2.8) rectangle (\xa+0.12,-5.8);
% AI Engine activo durante la atribución, mucho más larga
\filldraw[draw=turquesa!70, fill=turquesa!20]
 (\xa-0.12,-8.0) rectangle (\xa+0.12,-10.8);

%% ---- MENSAJES ----

% 1. Usuario → Frontend
\draw[sync] (\xu,-1.0) -- node[msglbl]{1:\enspace envía informe clínico} (\xf,-1.0);

% 2. Frontend → Backend (WebSocket)
\draw[sync] (\xf,-1.8) -- node[msglbl]{2:\enspace [WS] analyze\_report} (\xb,-1.8);

% 3. Backend: despacha comando (auto-mensaje)
\draw[sync, rounded corners=4pt]
 (\xb,-2.4) -- ++(0.6,0) -- ++(0,-0.55) -- node[right, font=\scriptsize\sffamily]
 {3: dispatch(AnalyzeReport)} ++(-0.6,0);

% 4. Backend → AI Engine
\draw[sync] (\xb,-3.6) -- node[msglbl]{4:\enspace POST /predict + /summarize/stream (paralelo)} (\xa,-3.6);

% 5. AI Engine → Backend (respuesta)
\draw[ret] (\xa,-5.2) -- node[msglbl]{5:\enspace \{codes, confidence\_scores\}} (\xb,-5.2);

% 6. Backend: despacha evento (auto-mensaje)
\draw[sync, rounded corners=4pt]
 (\xb,-5.9) -- ++(0.6,0) -- ++(0,-0.5) -- node[right, font=\scriptsize\sffamily]
 {6: dispatch(ReceiveAIPrediction) → Event Store} ++(-0.6,0);

% 7. Backend → Frontend: tarjeta de códigos
\draw[ret] (\xb,-6.9) -- node[msglbl]{7:\enspace [WS] analysis\_card\_received: codes} (\xf,-6.9);

% 8. Backend → AI Engine: atribución en segunda llamada, no bloquea
\draw[sync] (\xb,-7.8) -- node[msglbl]{8:\enspace POST /explain (no bloquea)} (\xa,-7.8);

% 9. Backend → Frontend: resumen en streaming (token a token)
\draw[ret] (\xb,-8.7) -- node[msglbl]{9:\enspace [WS] summary\_token ($\times$N)} (\xf,-8.7);

% 10. Backend → Frontend: resumen completo
\draw[ret] (\xb,-9.5) -- node[msglbl]{10:\enspace [WS] card summary + analysis\_complete} (\xf,-9.5);

% 11. AI Engine → Backend: términos explicativos
\draw[ret] (\xa,-11.0) -- node[msglbl]{11:\enspace \{triggers, method\}} (\xb,-11.0);

% 12. Backend → Frontend: términos, cuando llegan
\draw[ret] (\xb,-11.8) -- node[msglbl]{12:\enspace [WS] triggers\_received} (\xf,-11.8);

% 13. Frontend → Usuario
\draw[ret] (\xf,-12.6) -- node[msglbl]{13:\enspace completa los términos de cada código} (\xu,-12.6);

%% ---- NOTAS LATERALES ----
\node[note, right=0.3cm] at (\xa,-4.4)
 {Inferencia BERT\\y resumen (LLM)};
\node[note, right=0.3cm] at (\xa,-9.4)
 {Atribución por\\enmascarado:\\dos órdenes de\\magnitud más lenta\\que la predicción};

\end{tikzpicture}
 dentro de figure+\centering+\resizebox{\textwidth}{!}{...}
%
% Participantes (x): Usuario=0, Frontend=3.6, Backend=7.2, AI Engine=11
% Tiempo (y): crece hacia abajo (negativo en tikz)

\begin{tikzpicture}[
 every node/.style={font=\small\sffamily},
 % Cabecera de participante
 phead/.style={
 draw=aquaESI!80, fill=aquaESI!20, thick,
 rectangle, rounded corners=3pt,
 minimum width=2.4cm, minimum height=0.75cm, align=center
 },
 % Línea de vida
 life/.style={draw=gray!40, thick, dashed},
 % Mensaje síncrono (flecha llena)
 sync/.style={->, >=Stealth, thick, aquaESI!80},
 % Mensaje de retorno (flecha discontinua)
 ret/.style={->, >=Stealth, dashed, thick, gray!60},
 % Nota lateral
 note/.style={
 draw=ocre!60, fill=ocre!10,
 rectangle, rounded corners=2pt,
 font=\scriptsize\sffamily, inner sep=3pt, align=left
 },
 % Etiqueta de mensaje
 msglbl/.style={font=\scriptsize\sffamily, fill=white,
 inner sep=1pt, above, midway}
]

%% ---- COORDENADAS X de cada participante ----
\def\xu{0} % Usuario
\def\xf{3.6} % Frontend
\def\xb{7.2} % Backend
\def\xa{11.0} % AI Engine

%% ---- CABECERAS ----
\node[phead] at (\xu, 0) {\textbf{Usuario}};
\node[phead] at (\xf, 0) {\textbf{Frontend}};
\node[phead] at (\xb, 0) {\textbf{Backend}};
\node[phead] at (\xa, 0) {\textbf{AI Engine}};

%% ---- LÍNEAS DE VIDA ----
\draw[life] (\xu,-0.38) -- (\xu,-13.0);
\draw[life] (\xf,-0.38) -- (\xf,-13.0);
\draw[life] (\xb,-0.38) -- (\xb,-13.0);
\draw[life] (\xa,-0.38) -- (\xa,-13.0);

%% ---- ACTIVACIONES (barras sobre línea de vida) ----
% Backend activo durante la llamada al AI Engine
\filldraw[draw=aquaESI!70, fill=aquaESI!25]
 (\xb-0.12,-2.2) rectangle (\xb+0.12,-11.6);
% AI Engine activo durante la predicción
\filldraw[draw=turquesa!70, fill=turquesa!20]
 (\xa-0.12,-2.8) rectangle (\xa+0.12,-5.8);
% AI Engine activo durante la atribución, mucho más larga
\filldraw[draw=turquesa!70, fill=turquesa!20]
 (\xa-0.12,-8.0) rectangle (\xa+0.12,-10.8);

%% ---- MENSAJES ----

% 1. Usuario → Frontend
\draw[sync] (\xu,-1.0) -- node[msglbl]{1:\enspace envía informe clínico} (\xf,-1.0);

% 2. Frontend → Backend (WebSocket)
\draw[sync] (\xf,-1.8) -- node[msglbl]{2:\enspace [WS] analyze\_report} (\xb,-1.8);

% 3. Backend: despacha comando (auto-mensaje)
\draw[sync, rounded corners=4pt]
 (\xb,-2.4) -- ++(0.6,0) -- ++(0,-0.55) -- node[right, font=\scriptsize\sffamily]
 {3: dispatch(AnalyzeReport)} ++(-0.6,0);

% 4. Backend → AI Engine
\draw[sync] (\xb,-3.6) -- node[msglbl]{4:\enspace POST /predict + /summarize/stream (paralelo)} (\xa,-3.6);

% 5. AI Engine → Backend (respuesta)
\draw[ret] (\xa,-5.2) -- node[msglbl]{5:\enspace \{codes, confidence\_scores\}} (\xb,-5.2);

% 6. Backend: despacha evento (auto-mensaje)
\draw[sync, rounded corners=4pt]
 (\xb,-5.9) -- ++(0.6,0) -- ++(0,-0.5) -- node[right, font=\scriptsize\sffamily]
 {6: dispatch(ReceiveAIPrediction) → Event Store} ++(-0.6,0);

% 7. Backend → Frontend: tarjeta de códigos
\draw[ret] (\xb,-6.9) -- node[msglbl]{7:\enspace [WS] analysis\_card\_received: codes} (\xf,-6.9);

% 8. Backend → AI Engine: atribución en segunda llamada, no bloquea
\draw[sync] (\xb,-7.8) -- node[msglbl]{8:\enspace POST /explain (no bloquea)} (\xa,-7.8);

% 9. Backend → Frontend: resumen en streaming (token a token)
\draw[ret] (\xb,-8.7) -- node[msglbl]{9:\enspace [WS] summary\_token ($\times$N)} (\xf,-8.7);

% 10. Backend → Frontend: resumen completo
\draw[ret] (\xb,-9.5) -- node[msglbl]{10:\enspace [WS] card summary + analysis\_complete} (\xf,-9.5);

% 11. AI Engine → Backend: términos explicativos
\draw[ret] (\xa,-11.0) -- node[msglbl]{11:\enspace \{triggers, method\}} (\xb,-11.0);

% 12. Backend → Frontend: términos, cuando llegan
\draw[ret] (\xb,-11.8) -- node[msglbl]{12:\enspace [WS] triggers\_received} (\xf,-11.8);

% 13. Frontend → Usuario
\draw[ret] (\xf,-12.6) -- node[msglbl]{13:\enspace completa los términos de cada código} (\xu,-12.6);

%% ---- NOTAS LATERALES ----
\node[note, right=0.3cm] at (\xa,-4.4)
 {Inferencia BERT\\y resumen (LLM)};
\node[note, right=0.3cm] at (\xa,-9.4)
 {Atribución por\\enmascarado:\\dos órdenes de\\magnitud más lenta\\que la predicción};

\end{tikzpicture}
 dentro de figure+\centering+\resizebox{\textwidth}{!}{...}
%
% Participantes (x): Usuario=0, Frontend=3.6, Backend=7.2, AI Engine=11
% Tiempo (y): crece hacia abajo (negativo en tikz)

\begin{tikzpicture}[
 every node/.style={font=\small\sffamily},
 % Cabecera de participante
 phead/.style={
 draw=aquaESI!80, fill=aquaESI!20, thick,
 rectangle, rounded corners=3pt,
 minimum width=2.4cm, minimum height=0.75cm, align=center
 },
 % Línea de vida
 life/.style={draw=gray!40, thick, dashed},
 % Mensaje síncrono (flecha llena)
 sync/.style={->, >=Stealth, thick, aquaESI!80},
 % Mensaje de retorno (flecha discontinua)
 ret/.style={->, >=Stealth, dashed, thick, gray!60},
 % Nota lateral
 note/.style={
 draw=ocre!60, fill=ocre!10,
 rectangle, rounded corners=2pt,
 font=\scriptsize\sffamily, inner sep=3pt, align=left
 },
 % Etiqueta de mensaje
 msglbl/.style={font=\scriptsize\sffamily, fill=white,
 inner sep=1pt, above, midway}
]

%% ---- COORDENADAS X de cada participante ----
\def\xu{0} % Usuario
\def\xf{3.6} % Frontend
\def\xb{7.2} % Backend
\def\xa{11.0} % AI Engine

%% ---- CABECERAS ----
\node[phead] at (\xu, 0) {\textbf{Usuario}};
\node[phead] at (\xf, 0) {\textbf{Frontend}};
\node[phead] at (\xb, 0) {\textbf{Backend}};
\node[phead] at (\xa, 0) {\textbf{AI Engine}};

%% ---- LÍNEAS DE VIDA ----
\draw[life] (\xu,-0.38) -- (\xu,-13.0);
\draw[life] (\xf,-0.38) -- (\xf,-13.0);
\draw[life] (\xb,-0.38) -- (\xb,-13.0);
\draw[life] (\xa,-0.38) -- (\xa,-13.0);

%% ---- ACTIVACIONES (barras sobre línea de vida) ----
% Backend activo durante la llamada al AI Engine
\filldraw[draw=aquaESI!70, fill=aquaESI!25]
 (\xb-0.12,-2.2) rectangle (\xb+0.12,-11.6);
% AI Engine activo durante la predicción
\filldraw[draw=turquesa!70, fill=turquesa!20]
 (\xa-0.12,-2.8) rectangle (\xa+0.12,-5.8);
% AI Engine activo durante la atribución, mucho más larga
\filldraw[draw=turquesa!70, fill=turquesa!20]
 (\xa-0.12,-8.0) rectangle (\xa+0.12,-10.8);

%% ---- MENSAJES ----

% 1. Usuario → Frontend
\draw[sync] (\xu,-1.0) -- node[msglbl]{1:\enspace envía informe clínico} (\xf,-1.0);

% 2. Frontend → Backend (WebSocket)
\draw[sync] (\xf,-1.8) -- node[msglbl]{2:\enspace [WS] analyze\_report} (\xb,-1.8);

% 3. Backend: despacha comando (auto-mensaje)
\draw[sync, rounded corners=4pt]
 (\xb,-2.4) -- ++(0.6,0) -- ++(0,-0.55) -- node[right, font=\scriptsize\sffamily]
 {3: dispatch(AnalyzeReport)} ++(-0.6,0);

% 4. Backend → AI Engine
\draw[sync] (\xb,-3.6) -- node[msglbl]{4:\enspace POST /predict + /summarize/stream (paralelo)} (\xa,-3.6);

% 5. AI Engine → Backend (respuesta)
\draw[ret] (\xa,-5.2) -- node[msglbl]{5:\enspace \{codes, confidence\_scores\}} (\xb,-5.2);

% 6. Backend: despacha evento (auto-mensaje)
\draw[sync, rounded corners=4pt]
 (\xb,-5.9) -- ++(0.6,0) -- ++(0,-0.5) -- node[right, font=\scriptsize\sffamily]
 {6: dispatch(ReceiveAIPrediction) → Event Store} ++(-0.6,0);

% 7. Backend → Frontend: tarjeta de códigos
\draw[ret] (\xb,-6.9) -- node[msglbl]{7:\enspace [WS] analysis\_card\_received: codes} (\xf,-6.9);

% 8. Backend → AI Engine: atribución en segunda llamada, no bloquea
\draw[sync] (\xb,-7.8) -- node[msglbl]{8:\enspace POST /explain (no bloquea)} (\xa,-7.8);

% 9. Backend → Frontend: resumen en streaming (token a token)
\draw[ret] (\xb,-8.7) -- node[msglbl]{9:\enspace [WS] summary\_token ($\times$N)} (\xf,-8.7);

% 10. Backend → Frontend: resumen completo
\draw[ret] (\xb,-9.5) -- node[msglbl]{10:\enspace [WS] card summary + analysis\_complete} (\xf,-9.5);

% 11. AI Engine → Backend: términos explicativos
\draw[ret] (\xa,-11.0) -- node[msglbl]{11:\enspace \{triggers, method\}} (\xb,-11.0);

% 12. Backend → Frontend: términos, cuando llegan
\draw[ret] (\xb,-11.8) -- node[msglbl]{12:\enspace [WS] triggers\_received} (\xf,-11.8);

% 13. Frontend → Usuario
\draw[ret] (\xf,-12.6) -- node[msglbl]{13:\enspace completa los términos de cada código} (\xu,-12.6);

%% ---- NOTAS LATERALES ----
\node[note, right=0.3cm] at (\xa,-4.4)
 {Inferencia BERT\\y resumen (LLM)};
\node[note, right=0.3cm] at (\xa,-9.4)
 {Atribución por\\enmascarado:\\dos órdenes de\\magnitud más lenta\\que la predicción};

\end{tikzpicture}
}
\caption{Diagrama de secuencia del flujo de análisis clínico.}
\label{fig:seq-analisis}
\end{figure}

\section{API REST: frontend \texorpdfstring{$\rightarrow$}{->} backend}

La base de la URL es configurable mediante la variable de entorno \texttt{NEXT\_PUBLIC\_API\_URL} (valor por defecto: \texttt{http://localhost:4000/api}). Las rutas protegidas requieren la cabecera \texttt{Authorization: Bearer <token>}, donde \texttt{<token>} es un token firmado con \texttt{Phoenix.Token} emitido al iniciar sesión, con una validez de una hora.

La especificación completa de la API está disponible a través de la interfaz web de Swagger UI en \texttt{http://localhost:4000/docs}, donde es posible explorar todos los \emph{endpoints}, consultar los esquemas de petición y respuesta, y ejecutar llamadas de prueba directamente desde el navegador.

\subsection{Autenticación}

Los endpoints de registro, inicio de sesión, confirmación de email y cambio de idioma se listan en la Tabla~\ref{tab:api-auth}.

\begin{table}[h]
\centering
\caption{Endpoints de autenticación.}
\label{tab:api-auth}
\begin{tabular}{llp{7cm}}
\hline
\textbf{Método} & \textbf{Ruta} & \textbf{Descripción} \\
\hline
POST & \texttt{/auth/register} & Registra un nuevo usuario. Envía email de confirmación. \\
POST & \texttt{/auth/login} & Autentica al usuario; devuelve \texttt{\{user, token\}}. \\
GET & \texttt{/auth/confirm/:token} & Confirma el email del usuario a partir del enlace recibido. \\
PUT & \texttt{/users/locale} & Actualiza el idioma preferido del usuario. \emph{Protegido.} \\
\hline
\end{tabular}
\end{table}

El token se almacena en \texttt{localStorage} bajo la clave \texttt{cie10\_auth} como objeto \texttt{\{user, token\}}.

\subsection{Gestión de conversaciones}

El detalle de los endpoints para listar, archivar, restaurar y purgar conversaciones está en la Tabla~\ref{tab:api-conv}.

\begin{table}[h]
\centering
\caption{Endpoints de conversaciones.}
\label{tab:api-conv}
\begin{tabular}{lp{5cm}p{5cm}}
\hline
\textbf{Método} & \textbf{Ruta} & \textbf{Descripción} \\
\hline
GET & \texttt{/conversations?\allowbreak{}user\_id=<uuid>} & Lista conversaciones activas del usuario. \emph{Protegido.} \\
GET & \texttt{/conversations/trash?\allowbreak{}user\_id=<uuid>} & Lista conversaciones en papelera. \emph{Protegido.} \\
DELETE & \texttt{/conversations/:id} & Borrado lógico (soft delete). \emph{Protegido.} \\
PUT & \texttt{/conversations/\allowbreak{}:id/restore} & Restaura conversación de la papelera. \emph{Protegido.} \\
DELETE & \texttt{/conversations/\allowbreak{}:id/purge} & Borrado permanente e irreversible (sección~\ref{sec:rgpd}). \emph{Protegido.} \\
\hline
\end{tabular}
\end{table}

\subsection{Catálogo CIE-10-ES}

Los endpoints de búsqueda y consulta del catálogo se resumen en la Tabla~\ref{tab:api-cie10}.

\begin{table}[h]
\centering
\caption{Endpoints del catálogo CIE-10-ES (públicos, solo lectura).}
\label{tab:api-cie10}
\begin{tabular}{lp{5cm}p{5cm}}
\hline
\textbf{Método} & \textbf{Ruta} & \textbf{Descripción} \\
\hline
GET & \texttt{/cie10/search?\allowbreak{}q=<query>\&\allowbreak{}limit=<n>\&\allowbreak{}type=<tipo>} & Búsqueda en texto libre. \texttt{type} puede ser \texttt{diagnosis}, \texttt{procedure} o \texttt{chemical}. \\
GET & \texttt{/cie10/codes/:code} & Devuelve un código concreto con su descripción y metadatos. \\
GET & \texttt{/cie10/codes/\allowbreak{}:code/children} & Devuelve los códigos hijos en la jerarquía CIE-10. \\
\hline
\end{tabular}
\end{table}

\subsection{Traducciones}

El único endpoint de internacionalización aparece en la Tabla~\ref{tab:api-i18n}.

\begin{table}[h]
\centering
\caption{Endpoints de internacionalización (públicos).}
\label{tab:api-i18n}
\begin{tabular}{llp{7cm}}
\hline
\textbf{Método} & \textbf{Ruta} & \textbf{Descripción} \\
\hline
GET & \texttt{/translations/:locale} & Devuelve el mapa de cadenas de texto para \texttt{es} o \texttt{en}. \\
\hline
\end{tabular}
\end{table}

\section{Canal WebSocket: frontend \texorpdfstring{$\rightarrow$}{->} backend}

La conexión WebSocket se establece contra \texttt{ws://localhost:4000/socket} (configurable con \texttt{NEXT\_PUBLIC\_WS\_URL}) usando la librería \texttt{phoenix\_socket}. El token se pasa como parámetro de conexión para autenticar la sesión.

Cada conversación se suscribe a un canal propio con el topic \texttt{conversation:<conversation\_id>}, donde \texttt{conversation\_id} es un UUID generado por el cliente.

\subsection{Mensajes entrantes (cliente \texorpdfstring{$\rightarrow$}{->} servidor)}

El frontend envía siete tipos de mensajes al canal (Tabla~\ref{tab:ws-in}). El mensaje \texttt{join} se emite automáticamente al montar el componente de conversación y recibe como respuesta el historial completo del canal. Los mensajes de acción (\texttt{validate\_code}, \texttt{reject\_code} y \texttt{suggest\_code}) se generan cuando el codificador interactúa con los códigos predichos; cada uno emite el comando CQRS correspondiente sobre el agregado de análisis.

\begin{table}[H]
\centering
\caption{Mensajes enviados por el frontend al canal WebSocket.}
\label{tab:ws-in}
\begin{tabular}{lp{4.5cm}p{5.5cm}}
\hline
\textbf{Evento} & \textbf{Datos} & \textbf{Efecto} \\
\hline
\texttt{join} &  n/d & Une al cliente al canal; devuelve el historial de la conversación. \\
\texttt{send\_message} & \texttt{\{content: string\}} & Guarda el mensaje y lo difunde al canal. \\
\texttt{analyze\_report}& \texttt{\{report\_text: string\}} & Persiste el informe y lanza la tarea asíncrona de análisis IA. \\
\texttt{validate\_code} & \texttt{\{code\_id, cie10\_code\}} & Emite el comando \texttt{ValidateCode} al agregado. \\
\texttt{reject\_code} & \texttt{\{code\_id, cie10\_code, reason\}} & Emite el comando \texttt{RejectCode} al agregado. \\
\texttt{suggest\_code} & \texttt{\{selected\_text, suggested\_code\}} & Almacena una sugerencia manual del usuario. \\
\texttt{verify\_trigger} & \texttt{\{code\_id, trigger, verified\}} & Marca como verificado un término disparador. \\
\hline
\end{tabular}
\end{table}

\subsection{Mensajes salientes (servidor \texorpdfstring{$\rightarrow$}{->} cliente)}

Los mensajes que el backend emite hacia el cliente se listan en la Tabla~\ref{tab:ws-out}.

\begin{table}[h]
\centering
\caption{Mensajes emitidos por el backend hacia el cliente WebSocket.}
\label{tab:ws-out}
\begin{tabular}{lp{4.5cm}p{5cm}}
\hline
\textbf{Evento} & \textbf{Datos principales} & \textbf{Cuándo se emite} \\
\hline
\texttt{joined} & \texttt{\{history\}} & Al unirse al canal con éxito. \\
\texttt{new\_message} & \texttt{\{message\_id, content, user\_id, timestamp\}} & Cuando cualquier participante envía un mensaje. \\
\texttt{analysis\_started} & \texttt{\{message\_id\}} & Inmediatamente tras recibir \texttt{analyze\_report}. \\
\texttt{analysis\_card\_received}& \texttt{\{card\_id, card\_type, content, message\_id\}} & Por cada tarjeta generada (códigos y resumen). \\
\texttt{analysis\_complete} & \texttt{\{message\_id, predicted\_codes, engine\}} & Al completarse la predicción IA. \\
\texttt{code\_validated} & \texttt{\{code\_id, cie10\_code\}} & Tras confirmar la validación en el agregado. \\
\texttt{code\_rejected} & \texttt{\{code\_id, cie10\_code, reason\}} & Tras confirmar el rechazo en el agregado. \\
\texttt{summary\_token} & \texttt{\{message\_id, token\}} & Por cada fragmento del resumen generado en \emph{streaming}. \\
\texttt{triggers\_received} & \texttt{\{message\_id, method, triggers\}} & Cuando llega la respuesta de \texttt{/explain}, posterior a la de códigos. \\
\texttt{trigger\_verified} & \texttt{\{code\_id, trigger, verified\}} & Tras verificar un término disparador. \\
\texttt{analysis\_failed} & \texttt{\{message\_id, error\}} & Si el motor de IA devuelve un error. \\
\hline
\end{tabular}
\end{table}

El análisis genera la tarjeta de códigos (\texttt{codes}) a partir de \texttt{POST /predict} y, en paralelo, la tarjeta de resumen (\texttt{summary}) mediante \texttt{POST /summarize/stream}; ambas se emiten como eventos \texttt{analysis\_card\_received} antes del evento \texttt{analysis\_complete}.

\section{HTTP interno: backend \texorpdfstring{$\rightarrow$}{->} motor de IA}

El backend se comunica con el motor de IA mediante HTTP usando la librería Elixir \texttt{Req}. La URL base es \texttt{http://ai\_engine:8000} dentro de la red Docker, o \texttt{http://localhost:8000} en desarrollo sin Docker.

La especificación de los \emph{endpoints} del motor de IA está disponible a través de la interfaz web de Swagger UI en \texttt{http://localhost:8000/docs}, donde también se puede consultar el esquema OpenAPI generado automáticamente por FastAPI y realizar llamadas de prueba.

\subsection{Endpoint de predicción}

\begin{description}
 \item[Ruta] \texttt{POST /predict}
 \item[Cuerpo de la petición]~
\end{description}

\begin{verbatim}
{
 "text": "<informe clínico en español>",
 "engine": "bert" | "dict" | "both" | "fused", // opcional; por defecto "bert"
 "include_triggers": false                     // opcional; por defecto false
}
\end{verbatim}

\begin{description}
 \item[Respuesta] Objeto JSON con el campo \texttt{cards}, que contiene una única tarjeta de tipo \texttt{codes} con la lista de códigos predichos (cada uno con su descripción, confianza y términos que lo motivan). El resumen del informe se genera de forma independiente mediante \texttt{POST /summarize/stream}. El esquema completo de la petición y la respuesta se detalla en el Anexo~\ref{cap:AnexoM}.
\end{description}

La llamada se realiza dentro de una \texttt{Task.start} asíncrona, por lo que el backend no bloquea el canal WebSocket durante la inferencia. Cuando la respuesta llega, el backend emite el comando \texttt{ReceiveAIPrediction}, que genera los eventos correspondientes y activa los proyectores que persisten las tarjetas y los códigos predichos.

\subsection{Endpoint de explicabilidad}

\begin{description}
 \item[Ruta] \texttt{POST /explain}
 \item[Cuerpo] \texttt{\{text, codes, method?, top\_k?\}}
 \item[Respuesta] \texttt{\{method, triggers, unknown\_codes, timing\}}
\end{description}

La atribución de términos viaja en una llamada propia porque cuesta dos órdenes de magnitud más que la predicción. El parámetro \texttt{include\_triggers} de \texttt{/predict} está desactivado por defecto precisamente por eso: los códigos se devuelven en cuanto están y el backend lanza esta segunda petición en paralelo, emitiendo el evento \texttt{triggers\_received} por el canal WebSocket cuando la respuesta llega. La interfaz muestra entre tanto los términos que ya tiene (los del diccionario, si el motor es fusionado) junto a un indicador de progreso.

Si esta segunda llamada falla, la predicción no se invalida: se registra el aviso y la tarjeta conserva sus códigos sin términos explicativos.

\subsection{Endpoint de estado}

\begin{description}
 \item[Ruta] \texttt{GET /}
 \item[Respuesta] \texttt{\{status: \textquotedbl{}online\textquotedbl{}, model\_loaded: bool, dict\_loaded: bool, summarizer\_loaded: bool\}}
\end{description}

Permite al backend verificar que el motor de IA está disponible y que el modelo ha terminado de cargarse en memoria antes de aceptar peticiones de análisis.

\section{Mecanismo de resiliencia}

Si la llamada al motor de IA tiene éxito pero el despacho del comando \texttt{ReceiveAIPrediction} mediante Commanded falla, el backend dispone de un mecanismo de escritura directa en base de datos (\texttt{persist\_cards\_direct/3} y \texttt{persist\_predicted\_codes\_direct/2}) que garantiza que las tarjetas y los códigos predichos se guardan igualmente. Esto asegura que el usuario recibe una respuesta aunque el \emph{event store} experimente un problema transitorio.
